% Bagian: turing-machines
% Bab: machines-computations
% Subbagian: unary-numbers

\documentclass[../../../include/open-logic-section]{subfiles}

\begin{document}

\olfileid[id]{tur}{mac}{una}
\olsection{Representasi Uner Bilangan}

\begin{explain}
Mesin Turing bekerja pada barisan simbol yang tertulis pada pita. Bergantung
pada alfabet yang digunakan suatu mesin Turing, barisan simbol ini dapat
merepresentasikan beragam masukan dan keluaran. Tentu saja, yang paling
menarik ialah mesin Turing yang menghitung fungsi-fungsi \emph{aritmetis},
yakni fungsi bilangan asli. Cara sederhana merepresentasikan bilangan asli
ialah mengodekannya sebagai barisan satu simbol saja, yakni
$\TMstroke$. Jika $n \in \Nat$, misalkan $\TMstroke^n$ adalah barisan kosong
jika $n = 0$, dan selain itu merupakan barisan yang tepat terdiri atas $n$
simbol~$\TMstroke$.
\end{explain}

\begin{defn}[Komputasi]
Suatu mesin Turing~$M$ \emph{menghitung} fungsi $f\colon \Nat^k \to \Nat$ jika
dan hanya jika $M$~berhenti pada masukan
\[
\TMstroke^{n_1} \TMblank \TMstroke^{n_2} \TMblank \dots \TMblank \TMstroke^{n_k}
\]
dengan keluaran $\TMstroke^{f(n_1, \dots, n_k)}$.
\end{defn}

\begin{prob}
  Berikan definisi kapan suatu mesin Turing~$M$ menghitung fungsi
  $f\colon \Nat^k \to \Nat^m$.
\end{prob}

\begin{ex}\ollabel{ex:adder}
\emph{Penjumlahan:}
Marilah kita membangun mesin yang menghitung fungsi $f(n,m) = n + m$.
Untuk itu diperlukan mesin yang bermula dengan dua blok simbol~$\TMstroke$
sepanjang $n$ dan~$m$ pada pita, lalu berhenti dengan satu blok yang terdiri
atas $n+m$ simbol~$\TMstroke$. Kedua blok masukan simbol~$\TMstroke$
dipisahkan oleh sebuah~$\TMblank$, sehingga salah satu caranya ialah menulis
sebuah goresan pada petak yang memuat~$\TMblank$, lalu menghapus simbol
$\TMstroke$ terakhir.
\begin{figure}\[
\begin{tikzpicture}[->,>=stealth',shorten >=1pt,auto,node distance=2.8cm,
                    semithick]
  \tikzstyle{every state}=[fill=none,draw=black,text=black]

  \node[initial,state]         (A)              {$q_0$};
  \node[state]         (B) [right of=A] {$q_1$};
  \node[state]         (C) [right of=B] {$q_2$};

  \path (A) edge node {\TMtrans{\TMblank}{\TMstroke}{\TMstay}} (B)
           edge [loop above] node {\TMtrans{\TMstroke}{\TMstroke}{\TMright}} (A)
        (B) edge [loop above] node {\TMtrans{\TMstroke}{\TMstroke}{\TMright}} (B)
            edge node {\TMtrans{\TMblank}{\TMblank}{\TMleft}} (C)
        (C) edge [loop above] node {\TMtrans{\TMstroke}{\TMblank}{\TMstay}} (C);
\end{tikzpicture}
\]\caption{Mesin yang menghitung $f(x,y) = x+y$}
\ollabel{fig:adder}
\end{figure}
\end{ex}

\begin{prob}
Telusuri konfigurasi-konfigurasi mesin pada
\olref[tur][mac][una]{ex:adder} untuk masukan~$\tuple{3,2}$. Apa yang terjadi
jika mesin menghitung $0+0$?
\end{prob}

\begin{explain}
Pada \olref[rep]{ex:doubler}, kita memberikan contoh mesin Turing yang
menerima masukan berupa barisan simbol~$\TMstroke$ dan berhenti dengan barisan
yang berisi dua kali lebih banyak simbol~$\TMstroke$ pada pita---mesin
pengganda. Namun, karena keluarannya memuat simbol-simbol~$\TMblank$ di sebelah
kiri blok simbol~$\TMstroke$ yang telah digandakan, mesin itu sebenarnya tidak
menghitung fungsi $f(x) = 2x$, berlawanan dengan dugaan yang mungkin timbul.
Kita akan menguraikan dua cara untuk memperbaikinya.
\end{explain}

\begin{ex}
Mesin pada \olref{fig:doubler-disc} menghitung fungsi $f(x) = 2x$. Alih-alih
menghapus masukan dan menulis dua simbol~$\TMstroke$ di ujung kanan untuk
setiap simbol~$\TMstroke$ pada masukan seperti yang dilakukan mesin dari
\olref[rep]{ex:doubler}, mesin ini menambahkan satu simbol~$\TMstroke$ di
sebelah kanan untuk setiap simbol~$\TMstroke$ pada masukan. Mesin perlu
melacak tempat berakhirnya masukan, sehingga dibiarkannya sebuah~$\TMblank$ di
antara masukan dan goresan-goresan tambahan; pada bagian paling akhir, petak
itu diisinya dengan~$\TMstroke$. Kita juga harus ``mengingat'' posisi kita
dalam masukan, sehingga sebuah~$\TMstroke$ dalam blok masukan diganti sementara
dengan~$\TMblank$.
  \begin{figure}
\[
\begin{tikzpicture}[->,>=stealth',shorten >=1pt,auto,node distance=2.8cm,
                    semithick]
  \tikzstyle{every state}=[fill=none,draw=black,text=black]

  \node[initial,state] (0)              {$q_0$};
  \node[state]         (1) [above of=0] {$q_1$};
  \node[state]         (2) [above of=1] {$q_2$};
  \node[state]         (3) [right of=2] {$q_3$};
  \node[state]         (4) [below of=3] {$q_4$};
  \node[state]         (5) [below of=4] {$q_5$};
  \node[state]         (6) [above right of=4] {$q_6$};
  \node[state]         (7) [below of=6] {$q_7$};
  \node[state]         (8) [below of=7] {$q_8$};

  \path (0) edge node {\TMtrans{\TMstroke}{\TMblank}{\TMright}} (1)
%    (0) edge node {\TMtrans{\TMblank}{\TMblank}{\TMstay}} (h)
    (1) edge [loop left] node {\TMtrans{\TMstroke}{\TMstroke}{\TMright}} (1)
      edge node {\TMtrans{\TMblank}{\TMblank}{\TMright}} (2)
    (2) edge [loop above] node {\TMtrans{\TMstroke}{\TMstroke}{\TMright}} (2)
        edge node {\TMtrans{\TMblank}{\TMstroke}{\TMleft}} (3)
    (3) edge [loop above] node {\TMtrans{\TMstroke}{\TMstroke}{\TMleft}} (3)
        edge node[left] {\TMtrans{\TMblank}{\TMblank}{\TMleft}} (4)
    (4) edge        node[left] {\TMtrans{\TMstroke}{\TMstroke}{\TMleft}} (5)
    (5) edge [loop below] node {\TMtrans{\TMstroke}{\TMstroke}{\TMleft}} (5)
        edge              node {\TMtrans{\TMblank}{\TMstroke}{\TMright}} (0)
    (4) edge      node[sloped] {\TMtrans{\TMblank}{\TMstroke}{\TMright}} (6)
    (6) edge              node {\TMtrans{\TMblank}{\TMstroke}{\TMright}} (7)
    (7) edge [loop right] node {\TMtrans{\TMstroke}{\TMstroke}{\TMright}} (7)
        edge              node {\TMtrans{\TMblank}{\TMblank}{\TMleft}} (8)
    (8) edge [loop right] node {\TMtrans{\TMstroke}{\TMblank}{\TMstay}} (8);
    \end{tikzpicture}
\]
\caption{Mesin yang menghitung $f(x) = 2x$}
\ollabel{fig:doubler-disc}
\end{figure}
\end{ex}

\begin{ex}\ollabel{ex:mover}
Kemungkinan kedua untuk menghitung $f(x) = 2x$ ialah mempertahankan mesin
pengganda semula, tetapi menambahkan keadaan dan instruksi pada bagian akhirnya
untuk memindahkan blok goresan yang telah digandakan ke ujung kiri pita. Mesin
pada \olref{fig:mover} menjalankan tepat bagian terakhir ini: jika dijalankan
pada pita yang terdiri atas sebuah blok simbol~$\TMblank$, diikuti sebuah blok
simbol~$\TMstroke$ (dengan kepala berada di posisi mana pun dalam blok
$\TMblank$), mesin menghapus simbol-simbol~$\TMstroke$ satu demi satu dan
menuliskannya di awal pita. Agar dapat mengetahui kapan tugasnya selesai,
mesin terlebih dahulu menandai akhir blok simbol~$\TMstroke$ dengan simbol
$\TMendtape$, yang dihapus pada akhirnya. Kita memulai penomoran keadaan
dari~$q_6$ agar semua keadaan itu dapat ditambahkan ke mesin pengganda. Yang
masih diperlukan hanyalah instruksi tambahan $\delta(q_0,
\TMblank) = \tuple{q_6,\TMblank,\TMstay}$, yakni panah dari~$q_0$ ke~$q_6$
berlabel~$\TMtrans{\TMblank}{\TMblank}{\TMstay}$. (Ada satu masalah halus:
mesin yang dihasilkan tidak bekerja untuk masukan~$x=0$. Kita serahkan masalah
ini sebagai latihan.)
\begin{figure}
  \[
  \begin{tikzpicture}[->,>=stealth',shorten >=1pt,auto,node distance=2.8cm,
                      semithick]
    \tikzstyle{every state}=[fill=none,draw=black,text=black]
    \node[initial,state] (6)              {$q_6$};
    \node[state]         (7) [right of=6] {$q_7$};
    \node[state]         (8) [right of=7] {$q_8$};
    \node[state]         (9) [below of=8] {$q_9$};
    \node[state]         (10) [left of=9] {$q_{10}$};
    \node[state]         (11) [left of=10]  {$q_{11}$};
    \node[state]         (12) [below of=11]  {$q_{12}$};
    \node[state]         (13) [right of=12] {$q_{13}$};
    \node[state]         (14) [right of=13] {$q_{14}$};
    \path
    (6)  edge node {\TMtrans{\TMblank}{\TMblank}{\TMright}} (7)
    (7)  edge [loop above] node {\TMtrans{\TMblank}{\TMblank}{\TMright}} (7)
         edge node {\TMtrans{\TMstroke}{\TMstroke}{\TMright}} (8)
    (8)  edge [loop above] node {\TMtrans{\TMstroke}{\TMstroke}{\TMright}} (8)
         edge node {\TMtrans{\TMblank}{\TMendtape}{\TMleft}} (9)
    (9)  edge [loop right] node {\TMtrans{\TMstroke}{\TMstroke}{\TMleft}} (9)
         edge node {\TMtrans{\TMblank}{\TMblank}{\TMright}} (10)
    (10) edge node[above] {\TMtrans{\TMstroke}{\TMblank}{\TMleft}} (11)
    (11) edge [loop above] node {\TMtrans{\TMblank}{\TMblank}{\TMleft}} (11)
         edge node[left] {\begin{tabular}{@{}l@{}}
          \TMtrans{\TMendtape}{\TMendtape}{\TMright}\\
          \TMtrans{\TMstroke}{\TMstroke}{\TMright}
         \end{tabular}} (12)
    (12) edge node[] {\TMtrans{\TMblank}{\TMstroke}{\TMright}} (13)
    (13) edge [loop below] node {\TMtrans{\TMblank}{\TMblank}{\TMright}} (13)
         edge node[sloped] {\TMtrans{\TMstroke}{\TMblank}{\TMleft}} (11)
    (13) edge node {\TMtrans{\TMendtape}{\TMblank}{\TMstay}} (14);
  \end{tikzpicture}
  \]
  \caption{Memindahkan sebuah blok simbol~$\TMstroke$ ke kiri}
  \ollabel{fig:mover}
\end{figure}
\end{ex}

\begin{prob}
Dalam \olref[tur][mac][una]{ex:mover}, kita menguraikan sebuah mesin yang
terdiri atas gabungan mesin pengganda dari
\olref[tur][mac][una]{fig:doubler-disc} dan mesin pemindah dari
\olref[tur][mac][una]{fig:mover}. Apa yang terjadi jika mesin gabungan ini
dijalankan pada masukan~$x=0$, yakni pada pita kosong? Bagaimana Anda akan
memperbaiki mesin tersebut agar dalam hal ini mesin berhenti dengan
keluaran~$2x=0$? (Anda seharusnya dapat melakukannya dengan menambahkan satu
keadaan dan satu transisi.)
\end{prob}

\begin{prob}
\emph{Pengurangan:} Rancang mesin Turing yang, ketika diberi masukan berupa dua
untai goresan tak-kosong sepanjang $n$ dan~$m$, dengan $n > m$, menghitung
fungsi $f(n,m) = n - m$.
\end{prob}

\begin{prob}
\emph{Kesamaan:} Rancang mesin Turing untuk menghitung fungsi berikut:
\[
\fn{equality}(n,m) =
\begin{cases}
  \text{1} & \text{jika~$n = m$} \\
  \text{0} & \text{jika~$n \neq m$}
\end{cases}
\]
dengan~$n$ dan~$m \in \PosInt$.
\end{prob}

\begin{prob}
Rancang mesin Turing untuk menghitung fungsi $\min(x,y)$, dengan $x$ dan $y$
bilangan bulat positif yang direpresentasikan pada pita oleh untai-untai
simbol~$\TMstroke$ yang dipisahkan oleh sebuah~$\TMblank$. Anda boleh
menggunakan simbol tambahan dalam alfabet mesin.

Fungsi $\min$ memilih nilai terkecil di antara argumen-argumennya, sehingga
$\min(3,5)=3$, $\min(20,16)=16$, dan $\min(4,4)=4$, dan seterusnya.
\end{prob}

\begin{defn}
  Suatu mesin Turing~$M$ menghitung fungsi parsial $f\colon \Nat^k
  \pto \Nat$ jika dan hanya jika
  \begin{enumerate}
    \item $M$ berhenti pada masukan $\TMstroke^{n_1}\concat\TMblank\concat
    \dots \concat\TMblank\concat\TMstroke^{n_k}$ dengan keluaran
    $\TMstroke^{m}$ jika $f(n_1, \dots, n_k) = m$,
    \item $M$ sama sekali tidak berhenti, atau berhenti dengan keluaran yang
    bukan satu blok tunggal simbol~$\TMstroke$, jika $f(n_1, \dots, n_k)$ tidak
    terdefinisi.
  \end{enumerate}
\end{defn}

\end{document}
